%! Matrix Computations and A = LU — Unit 2, Lesson 2.3 — Homework
\documentclass[10pt]{article}
\usepackage{linalg-article}
\usepackage{linalg-boxes}
\newcommand{\TallMath}[1]{$\displaystyle #1\rule[-1.4em]{0pt}{3.2em}$}
\newcommand{\xx}{\mathbf{x}}
\newcommand{\bb}{\mathbf{b}}
\newcommand{\cc}{\mathbf{c}}

\begin{document}

\pageheader{Unit 2, Lesson 2.3}{Homework}
\namedateperiod

\vspace{0.12in}

\begin{notesbox}{Practice}
Show your steps. Remember: $U$ is elimination's result, $L$ holds the multipliers ($+\ell$) with $1$s
on the diagonal, and $A\xx=\bb$ solves in two sweeps --- $L\cc=\bb$ forward, then $U\xx=\cc$ back.

\begin{enumerate}[leftmargin=1.9em, itemsep=12pt]
  \item \textbf{Factor.} For each $A$, name the pivot and multiplier $\ell$, then write $L$ and $U$.
    \begin{enumerate}[label=(\alph*), leftmargin=1.8em, itemsep=4pt]
      \item $A=\begin{bmatrix} 1 & 5 \\ 2 & 12 \end{bmatrix}$ \hfill
      \item $A=\begin{bmatrix} 2 & 3 \\ 8 & 15 \end{bmatrix}$
    \end{enumerate}
    \vspace{1.5cm}
  \item \textbf{Check.} For both matrices in Problem~1, multiply $LU$ and confirm you get back $A$.
    \vspace{1.5cm}
  \item \textbf{Two sweeps.} Factor $A=\begin{bmatrix} 1 & 2 \\ 4 & 9 \end{bmatrix}$, then solve
        $A\xx=\bb$ for $\bb=\begin{bmatrix} 3 \\ 14 \end{bmatrix}$ using $L\cc=\bb$ then $U\xx=\cc$.
    \vspace{1.8cm}
  \item \textbf{Juice bar --- one factorization, two batches.} Each liter of \textbf{Mix P} carries $1$
        unit sugar and $2$ units vitamin; each liter of \textbf{Mix Q} carries $2$ sugar and $6$
        vitamin. With $p$ liters of P and $q$ of Q the totals are $A\begin{bsmallmatrix} p\\q \end{bsmallmatrix}=\bb$,
        $A=\begin{bmatrix} 1 & 2 \\ 2 & 6 \end{bmatrix}$ (sugar row, vitamin row).
    \begin{enumerate}[label=(\alph*), leftmargin=1.8em, itemsep=6pt]
      \item Factor $A=LU$ once. \vspace{1.0cm}
      \item \textbf{Batch 1:} sugar $5$, vitamin $12$. Two sweeps --- find $p,q$. \vspace{1.0cm}
      \item \textbf{Batch 2:} sugar $8$, vitamin $22$. Same $L,U$ --- find $p,q$. \vspace{1.0cm}
      \item \emph{Interpret:} give the liters for each batch, and say why one factorization handled both. \writeline
    \end{enumerate}
  \item \textbf{A $3\times3$.} Factor $A=\begin{bmatrix} 1 & 1 & 2 \\ 3 & 5 & 7 \\ 1 & 5 & 5 \end{bmatrix}$
        (collect $\ell_{21},\ell_{31},\ell_{32}$ into $L$), then solve $A\xx=\bb$ for
        $\bb=\begin{bmatrix} 5 \\ 20 \\ 16 \end{bmatrix}$ by two sweeps.
    \vspace{2.2cm}
  \item \textbf{Explain.} In Problem~3 the middle vector $\cc$ from $L\cc=\bb$ is exactly the
        right-hand side you would get by running elimination on $\bb$. Explain \emph{why} the forward
        sweep and ``updating $\bb$ during elimination'' give the same thing. \writeline
\end{enumerate}
\end{notesbox}

\begin{extensionbox}
\textbf{Where $L$ comes from.} In 2.2 you undid a sequence of steps with
$(E_{31}E_{21})^{-1}=E_{21}^{-1}E_{31}^{-1}$. Those undo-matrices, multiplied together, \emph{are} $L$:
\[
  L = E_{21}^{-1}E_{31}^{-1}E_{32}^{-1}
    = \begin{bmatrix} 1&0&0\\ \ell_{21}&1&0\\ \ell_{31}&\ell_{32}&1 \end{bmatrix}.
\]
Using the multipliers from Problem~5, write out this product and confirm the multipliers land in their
$L$-positions with no interference. Why is it so much cheaper to keep $L$ and $U$ than to build $A^{-1}$? \writeline
\end{extensionbox}

\begin{spiralbox}
\textbf{Coming next --- Lesson 2.4: Permutations and Transposes.} Every matrix here factored cleanly as
$A=LU$. But what if a pivot is $0$ and you must swap rows first? A \textbf{permutation matrix} $P$
records the swap, and the factorization becomes $PA=LU$.
\end{spiralbox}

\end{document}
